% Auto-generated by validate_t4_risk_dominance_sample_complexity.py
\begin{tabular}{llccccc}
\toprule
\textbf{Scenario} & $\pi$ & \textbf{Margin} & \textbf{Margin LCB} & \textbf{$n$ used} & \textbf{$n^*$ needed} & \textbf{Dominates?} \\
\midrule
zero\_attack\_k4 & 0.05 & -0.0002 & -0.0036 & 1599 & $\infty$ & no \\
zero\_attack\_k4 & 0.10 & -0.0004 & -0.0072 & 1599 & $\infty$ & no \\
zero\_attack\_k4 & 0.25 & -0.0010 & -0.0179 & 1599 & $\infty$ & no \\
zero\_attack\_k4 & 0.50 & -0.0019 & -0.0359 & 1599 & $\infty$ & no \\
max\_attack\_k4 & 0.05 & +0.0005 & -0.0029 & 1600 & 82427 & no \\
max\_attack\_k4 & 0.10 & +0.0009 & -0.0058 & 1600 & 82427 & no \\
max\_attack\_k4 & 0.25 & +0.0024 & -0.0146 & 1600 & 82427 & no \\
max\_attack\_k4 & 0.50 & +0.0047 & -0.0292 & 1600 & 82427 & no \\
\bottomrule
\end{tabular}

% Margin LCB: Hoeffding worst-case on Delta_0, Delta_1 at 95% confidence,
% range R=2. n* = C^2 ln(2/delta) / (2 M0^2) is the min fired-sample count
% to certify dominance at prevalence pi (infinite when point margin <= 0).